\chapter{\ploren{Cel i zakres pracy}{Aim and scope of the thesis}}

W tym rozdziale powinno znaleźć się już szczegółowe sformułowanie celu pracy oraz poruszanej w niej problematyki, której kontekst został nakreślony w rozdziale wprowadzającym. Możliwe są m.in. następujące cele: analiza istniejących metod i technik, eksperymentalne badania różnych rozwiązań, zaprojektowanie i wykonanie urządzenia lub/oraz programu komputerowego.

Należy tu pokrótce omówić zawartość kolejnych rozdziałów pracy. Skoncentrować się przede wszystkim na opisaniu własnego wkładu w rozwój dziedziny, będącej przedmiotem pracy. Praca dyplomowa ma na celu rozwiązanie lub analizę konkretnego problemu. Praca magisterska powinna ponadto posiadać element nowatorstwa – należy go tu wskazać. Na jej podstawie ocenia się zdolność dyplomanta do prawidłowego formułowania i analizy problemów, umiejętność wyciągania wniosków, zbierania materiałów, przeprowadzania badań oraz analizy wyników.

„Pusty obszar” jak poniżej może w pracy pojawić się wyłącznie na końcu głównego rozdziału. Nowy rozdział zaczynamy od nowej strony.
